\chapter{Theoretical Background}
\label{ch:theoretical_background}

\epigraph{The cell is a machine driven by energy.}{--- Albert Szent-Gy\"{o}rgyi}

\lettrine[lines=2]{T}{his chapter} establishes the theoretical foundations of the RCC agent-based model, covering agent-based modeling and Repast4Py, the immunology and metabolism of the tumor microenvironment, the treatment modalities implemented in the simulation, and the related work that positions this thesis within the computational oncology literature.

% ============================================================================
\section{Agent-Based Modeling}
\label{sec:bg_abm}
% ============================================================================

An agent-based model (ABM) consists of autonomous agents, a shared environment, and behavioral rules.  System-level behavior is emergent rather than prescribed by global equations, making ABMs powerful for systems with heterogeneity, spatial structure, and nonlinear feedback---hallmarks of the tumor microenvironment.  In contrast to ODE/PDE approaches that describe aggregate state variables, ABMs represent each entity individually (Figure~\ref{fig:abm_vs_ode}).  In this model, tumor cells carry individual 16-gene DNA sequences, agents must be in spatial proximity to interact, and mutation and immune recognition follow per-agent probabilistic rules.  This richness comes at a computational cost scaling with population size, motivating MPI-based parallelization via Repast4Py.

\begin{figure}[htbp]
\centering
\adjustbox{max width=0.85\textwidth, max height=0.7\textheight}{% ABM vs ODE/PDE Comparison Diagram
\begin{tikzpicture}[scale=0.8, every node/.style={scale=0.8}]

% Title
\node[text width=18cm, text centered] at (9,14) {
    \Large\textbf{Modeling Paradigms: Agent-Based vs Equation-Based Approaches}
};

% Divide into two sides
\draw[thick, dashed] (9,0) -- (9,13);

% ABM Side (Left)
\node[text width=8cm, text centered] at (4.5,12.5) {\textbf{\textcolor{immunecell}{Agent-Based Model (ABM)}}};

% Individual agents
\node[draw, circle, fill=tumorcell, text=white, minimum size=0.8cm] at (2,11) {TC};
\node[draw, circle, fill=immunecell, text=white, minimum size=0.6cm] at (3,10.5) {T};
\node[draw, circle, fill=immunecell, text=white, minimum size=0.6cm] at (1.5,10.5) {M};
\node[draw, circle, fill=immunecell, text=white, minimum size=0.6cm] at (2.5,10) {NK};
\node[draw, circle, fill=bloodvessel, text=white, minimum size=0.6cm] at (1,11.5) {BV};

% Grid representation
\draw[step=0.3cm, gray!30] (4,9.5) grid (7,11.5);
\node[above] at (5.5,11.5) {\textbf{3D Spatial Grid}};

% Individual cell properties
\node[draw, rounded corners, fill=tumorcell!20, text width=3.5cm] at (2,9) {
    \textbf{Individual Properties}\\
    • Unique DNA profile\\
    • Spatial position (x,y,z)\\
    • Activation state\\
    • Age/lifecycle\\
    • Local interactions
};

% Stochastic rules
\node[draw, rounded corners, fill=blue!20, text width=3.5cm] at (6,9) {
    \textbf{Stochastic Rules}\\
    • Probabilistic decisions\\
    • Random movement\\
    • Bernoulli events\\
    • Monte Carlo sampling
};

% Emergent behaviors
\node[draw, rounded corners, fill=green!20, text width=7cm] at (4.5,7.5) {
    \textbf{Emergent System Behaviors}\\
    • Tumor growth patterns from cell division rules\\
    • Immune response waves from local interactions\\
    • Spatial clustering without global coordination\\
    • Treatment resistance from mutation accumulation
};

% Advantages
\node[draw, rounded corners, fill=green!10, text width=7cm] at (4.5,6) {
    \textbf{ABM Advantages}\\
    + Natural heterogeneity representation\\
    + Spatial dynamics and local interactions\\
    + Stochasticity at individual level\\
    + Easy addition of new cell types\\
    + Direct biological interpretation\\
    + Modular/extensible architecture
};

% Disadvantages
\node[draw, rounded corners, fill=red!10, text width=7cm] at (4.5,4.5) {
    \textbf{ABM Disadvantages}\\
    - Computationally expensive\\
    - Many parameters to calibrate\\
    - Stochastic results require ensembles\\
    - Difficult mathematical analysis\\
    - Scale limitations (memory/time)
};

% Example equation
\node[draw, rounded corners, fill=gray!20, text width=7cm] at (4.5,3) {
    \textbf{Typical ABM Rules}\\
    $P_{kill} = f(\text{activation}, \text{glucose}, \text{effects})$\\
    $P_{divide} = g(\text{DNA}, \text{growth factors})$\\
    $\mathbf{x}_{t+1} = \mathbf{x}_t + \text{movement}(\text{targets})$\\
    IF random() $< P_{kill}$ THEN remove target
};

% ODE/PDE Side (Right)
\node[text width=8cm, text centered] at (13.5,12.5) {\textbf{\textcolor{purple}{Equation-Based Models (ODE/PDE)}}};

% Population-level variables
\node[draw, rounded corners, fill=purple!20, text width=3.5cm] at (11,11) {
    \textbf{Population Variables}\\
    $T(t)$ = Tumor density\\
    $I(t)$ = Immune density\\
    $V(t)$ = Vessel density\\
    $C(x,t)$ = Drug conc.
};

% Continuous functions
\node[draw, rounded corners, fill=orange!20, text width=3.5cm] at (15.5,11) {
    \textbf{Continuous Fields}\\
    Smooth functions\\
    Spatial averages\\
    Rate equations\\
    Analytical solutions
};

% Differential equations
\node[draw, rounded corners, fill=purple!30, text width=7cm] at (13.5,9.5) {
    \textbf{System of Differential Equations}\\
    $\frac{dT}{dt} = r_T T(1-T/K) - k_{TI} TI$\\[0.1cm]
    $\frac{dI}{dt} = s_I + \alpha TI - d_I I$\\[0.1cm]
    $\frac{\partial C}{\partial t} = D\nabla^2 C - \lambda C + S$
};

% Deterministic dynamics
\node[draw, rounded corners, fill=blue!30, text width=3.5cm] at (11,8) {
    \textbf{Deterministic}\\
    • Fixed trajectories\\
    • Reproducible results\\
    • Average behaviors\\
    • Smooth dynamics
};

% Mean-field approximation
\node[draw, rounded corners, fill=orange!30, text width=3.5cm] at (16,8) {
    \textbf{Mean-Field}\\
    • Population averages\\
    • Well-mixed assumption\\
    • Rate constants\\
    • Effective parameters
};

% Advantages
\node[draw, rounded corners, fill=green!10, text width=7cm] at (13.5,6) {
    \textbf{ODE/PDE Advantages}\\
    + Mathematically tractable\\
    + Fast computation\\
    + Analytical solutions possible\\
    + Parameter sensitivity analysis\\
    + Stability/bifurcation analysis\\
    + Well-established theory
};

% Disadvantages
\node[draw, rounded corners, fill=red!10, text width=7cm] at (13.5,4.5) {
    \textbf{ODE/PDE Disadvantages}\\
    - No individual heterogeneity\\
    - Limited spatial resolution\\
    - Difficult to add discrete events\\
    - Mean-field assumptions\\
    - Less intuitive for cell biology
};

% When to use each
\node[draw, rounded corners, fill=yellow!20, text width=7cm] at (4.5,1.5) {
    \textbf{Use ABM When:}\\
    • Individual cell behavior matters\\
    • Spatial patterns important\\
    • Discrete events (division, death)\\
    • Heterogeneity drives outcomes\\
    • Exploring mechanisms
};

\node[draw, rounded corners, fill=yellow!20, text width=7cm] at (13.5,1.5) {
    \textbf{Use ODE/PDE When:}\\
    • Population dynamics sufficient\\
    • Large-scale patterns\\
    • Parameter estimation needed\\
    • Mathematical analysis required\\
    • Fast screening/optimization
};

% Hybrid approaches
\node[draw, rounded corners, fill=cyan!20, text width=16cm, text centered] at (9,0.3) {
    \textbf{Hybrid Approaches:} ABM for cellular decisions + PDE for chemical fields (glucose) • Multi-scale modeling • ABM → ODE parameter estimation
};

% Arrows showing relationship
\draw[thick, <->, blue] (7.5,9) to[bend left=30] (10.5,9) node[midway, above] {\textbf{Parameter Mapping}};

% Complexity indicators
\node[draw, star, fill=red!50] at (1,0.8) {High Complexity};
\node[draw, star, fill=yellow!50] at (8,0.8) {Medium};
\node[draw, star, fill=green!50] at (16,0.8) {Low Complexity};

\end{tikzpicture}}
\caption{Comparison of agent-based and equation-based modeling paradigms. ABMs represent individual agents with local interaction rules, producing emergent system-level behavior; ODE/PDE models describe aggregate state variables through global equations.}
\label{fig:abm_vs_ode}
\end{figure}

\paragraph{Repast4Py Framework.}
\label{subsec:repast4py}

Repast4Py~\cite{collier2020} is a Python-based ABM toolkit that provides native MPI-based parallelism, distinguishing it from single-process frameworks such as Mesa or NetLogo. The simulation is distributed across MPI ranks, each responsible for a spatial partition and the agents within it. Agents are placed on a \texttt{SharedGrid}---a discrete $N$-dimensional lattice supporting multiple occupancy---and interact only with nearby agents, including read-only ghost copies of agents on neighboring ranks. This architecture enables near-linear scaling for spatially local interaction models like the tumor microenvironment. Chapter~\ref{ch:implementation} provides a detailed treatment of the framework's internals and their application in this work.

% ============================================================================
\section{Immunology of Renal Cell Carcinoma}
\label{sec:bg_immunology}
% ============================================================================

Renal cell carcinoma is one of the most immunogenic solid tumors, a property rooted in the frequent inactivation of the von Hippel--Lindau (VHL) tumor suppressor gene (70--80\% of clear cell cases)~\cite{gnarra1994}. VHL loss stabilizes hypoxia-inducible factors, driving VEGF overexpression and hypervascularization that facilitates immune cell infiltration. The resulting paradox---abundant tumor-infiltrating lymphocytes coexisting with progressive tumor growth---reflects the interplay between immune activation and tumor-mediated evasion that the model seeks to capture. High mutational burden and a diverse neoantigen repertoire further make RCC amenable to checkpoint immunotherapy~\cite{Motzer2015}. Table~\ref{tab:immune_cells} summarizes the thirteen immune cell types represented in the model, spanning both innate and adaptive arms of the anti-tumor response.

\begin{table}[htbp]
\centering
\caption{Immune cell types and primary functions in the RCC model.}
\label{tab:immune_cells}
{\footnotesize
\begin{tabular}{@{}ll@{}}
\toprule
\textbf{Cell Type} & \textbf{Primary Function} \\
\midrule
CD8+ Cytotoxic T Cell & Direct tumor killing \\
CD8+ / CD4+ Naive T Cell & CTL / helper precursors \\
Th1 / Th2 Helper T Cell & Pro-/anti-inflammatory \\
Regulatory T Cell & Immunosuppression \\
Dendritic Cell / pDC & Antigen presentation, T cell \& NK activation \\
M1 / M2 Macrophage & Anti-tumor killing / pro-tumor promotion \\
Natural Killer Cell & Innate tumor killing \\
Mast Cell / Neutrophil & Context-dependent / sentinel \\
\bottomrule
\end{tabular}
}
\end{table}

Among these, the PD-1/PD-L1 checkpoint axis is most clinically relevant: PD-L1 on tumor cells engages PD-1 on T cells, suppressing proliferation and cytotoxic activity---a mechanism directly targeted by the ICI treatment modeled in this work. Figure~\ref{fig:immune_response_cascade} illustrates the temporal progression from innate recognition through adaptive activation to effector-mediated killing, including the immunosuppressive feedback loops that can blunt the response.

\begin{figure}[htbp]
\centering
\adjustbox{max width=0.85\textwidth, max height=0.7\textheight}{% Immune Response Cascade — 4-Phase Horizontal Timeline
\begin{tikzpicture}[
    cell/.style={draw, circle, font=\footnotesize, text=white,
        minimum size=0.85cm, inner sep=0pt},
    cytokine/.style={draw, diamond, fill=glucose!70, font=\footnotesize,
        minimum size=0.55cm, inner sep=1pt},
    phaselabel/.style={font=\small\bfseries, anchor=south},
    arr/.style={->, >=stealth, thick},
    sarr/.style={->, >=stealth, thin, dashed}
]

% ── Phase separators (vertical dashed lines) ────────────────
\foreach \x in {3.4, 6.8, 10.2} {
    \draw[densely dashed, gray!45] (\x,-0.4) -- (\x,6.8);
}

% ── Phase labels ────────────────────────────────────────────
\node[phaselabel] at (1.7, 6.8) {Phase 1};
\node[font=\footnotesize, anchor=north] at (1.7,6.7) {Recognition};

\node[phaselabel] at (5.1, 6.8) {Phase 2};
\node[font=\footnotesize, anchor=north] at (5.1,6.7) {Activation};

\node[phaselabel] at (8.5, 6.8) {Phase 3};
\node[font=\footnotesize, anchor=north] at (8.5,6.7) {Effector};

\node[phaselabel] at (11.9, 6.8) {Phase 4};
\node[font=\footnotesize, anchor=north] at (11.9,6.7) {Resolution};

% ── Timeline arrow ──────────────────────────────────────────
\draw[arr, gray!60] (0,0) -- (13.6,0);
\node[font=\footnotesize, gray!70, anchor=west] at (13.7,0) {time};

% ══ Phase 1 — Innate Recognition ════════════════════════════
\node[cell, fill=bloodvessel]  (nk)    at (0.8, 4.5) {NK};
\node[cell, fill=bloodvessel]  (macro) at (1.7, 2.8) {M$\phi$};
\node[cell, fill=bloodvessel]  (dc)    at (2.6, 4.5) {DC};

\node[font=\footnotesize, gray, anchor=north] at (1.7,2.1) {innate cells};

% ══ Phase 2 — Activation ════════════════════════════════════
\node[cell, fill=immunecell]  (cd8) at (4.6, 4.5) {CD8$^+$};
\node[cell, fill=immunecell]  (cd4) at (5.6, 2.8) {CD4$^+$};

% DC → presents antigen
\draw[arr] (dc) -- (cd8)
    node[midway, above, font=\footnotesize] {Ag};
\draw[arr] (dc.east) ++(0.1,0) -- ++(0.5,0) |- (cd4.west)
    node[near start, above, font=\footnotesize] {};

% ══ Phase 3 — Effector ══════════════════════════════════════
\node[cell, fill=immunecell]  (ctl) at (7.8, 5.2) {CTL};
\node[cell, fill=immunecell]  (th1) at (8.2, 3.4) {Th1};
\node[cell, fill=immunecell]  (th2) at (8.8, 1.8) {Th2};

% Tumour target
\node[cell, fill=tumorcell, minimum size=1cm, font=\footnotesize]
    (tumour) at (10.0, 3.4) {TC};

% Activation → effector
\draw[arr] (cd8) -- (ctl);
\draw[arr] (cd4) -- (th1);
\draw[arr] (cd4.east) ++(0.1,0) -- ++(0.4,0) |- (th2.west);

% Killing arrows
\draw[arr, tumorcell]  (ctl) -- (tumour)
    node[midway, above, font=\footnotesize] {kill};
\draw[arr, tumorcell]  (nk) to[bend left=18] (tumour);
\draw[arr, immunecell, thin] (th1) -- (tumour)
    node[midway, below, font=\footnotesize] {support};

% Cytokine diamonds
\node[cytokine] (ifng) at (9.0, 5.6) {\scriptsize IFN-$\gamma$};
\node[cytokine] (il2)  at (9.6, 2.6) {\scriptsize IL-2};
\node[cytokine] (tnfa) at (7.4, 1.6) {\scriptsize TNF-$\alpha$};

\draw[sarr] (ctl) -- (ifng);
\draw[sarr] (th1) -- (il2);
\draw[sarr] (th1) -- (tnfa);

% ══ Phase 4 — Resolution ════════════════════════════════════
\node[cell, fill=immunecell!50] (tmem)  at (11.4, 5.2) {T$_M$};
\node[cell, fill=gray!60]       (treg)  at (12.2, 2.8) {Treg};

\draw[arr] (ctl) -- (tmem)
    node[midway, above, font=\footnotesize] {memory};
\draw[sarr, gray!70] (treg) to[bend right=20] (ctl)
    node[near start, below, font=\footnotesize] {suppress};

% ── Compact legend ──────────────────────────────────────────
\node[draw, rounded corners=2pt, fill=white, inner sep=4pt,
      font=\footnotesize, anchor=south west] at (0,-0.3) {%
    \textcolor{bloodvessel}{$\bullet$}\;Innate\quad
    \textcolor{immunecell}{$\bullet$}\;Adaptive\quad
    \textcolor{tumorcell}{$\bullet$}\;Tumor\quad
    \textcolor{gray}{$\bullet$}\;Suppressor\quad
    \textcolor{glucose!80!black}{$\diamond$}\;Cytokine};

\end{tikzpicture}
}
\caption{Immune response cascade from innate recognition through adaptive activation to effector-mediated tumor killing. Cytokines (IFN-$\gamma$, IL-2, TNF-$\alpha$) mediate cross-talk between phases. Regulatory T cells and M2 macrophages provide immunosuppressive feedback that can blunt the effector response.}
\label{fig:immune_response_cascade}
\end{figure}

The cascade can equivalently be viewed as a sequential process with regulatory feedback (Figure~\ref{fig:immune_cascade_bpmn}), where each step depends on the preceding one, making the response vulnerable to disruption at multiple points.

\begin{figure}[htbp]
\centering
\adjustbox{max width=0.85\textwidth, max height=0.7\textheight}{% Immune Cascade — BPMN-Style Flow Chart
\begin{tikzpicture}[
    action/.style={draw, rounded corners=3pt, fill=immunecell!10,
        font=\small, text width=2.8cm, align=center,
        minimum height=0.9cm, inner sep=4pt},
    actor/.style={draw, ellipse, fill=glucose!20,
        font=\small, minimum height=0.8cm, text width=2.0cm,
        align=center, inner sep=2pt},
    arr/.style={->, >=stealth, thick},
    fbarr/.style={->, >=stealth, thick, dashed}
]

% ── Main vertical flow ──────────────────────────────────────
\node[action, fill=tumorcell!15] (prolif) at (0, 0)
    {Tumor Cell\\Proliferation};

\node[actor] (dc) at (0,-2)
    {Dendritic\\Cell};

\node[action] (phago) at (0,-4)
    {Phagocytosis};

\node[action] (agpres) at (0,-6)
    {Antigen\\Presentation};

\node[action, fill=immunecell!15] (tact) at (0,-8)
    {T Cell\\Activation};

% ── Main cascade arrows ─────────────────────────────────────
\draw[arr] (prolif) -- (dc)
    node[midway, right, font=\footnotesize] {detection};
\draw[arr] (dc) -- (phago);
\draw[arr] (phago) -- (agpres);
\draw[arr] (agpres) -- (tact);

% ── CTL branch (right) ──────────────────────────────────────
\node[actor, fill=immunecell!15] (ctl) at (4.5,-8)
    {Cytotoxic\\T Cell};

\node[action, fill=tumorcell!15] (kill) at (4.5,-10)
    {Tumor Cell\\Killing};

\draw[arr] (tact) -- (ctl);
\draw[arr] (ctl) -- (kill);

% CTL feedback loop: killing → DC detection (dashed green)
\draw[fbarr, treatment] (kill.east) -- ++(1.2,0) |- (dc.east)
    node[pos=0.25, right, font=\footnotesize, treatment!70!black]
    {debris recycle};

% ── Treg branch (left) ──────────────────────────────────────
\node[actor, fill=gray!20] (treg) at (-4.5,-8)
    {Regulatory\\T Cell};

\node[action, fill=tumorcell!10] (suppr) at (-4.5,-10)
    {Immuno-\\suppression};

\draw[arr] (tact) -- (treg);
\draw[arr] (treg) -- (suppr);

% Treg feedback: suppression → T cell activation (dashed red)
\draw[fbarr, tumorcell] (suppr.east) -- ++(0.8,0) |- ([yshift=-3pt]tact.west)
    node[pos=0.25, right, font=\footnotesize, tumorcell!70!black]
    {inhibit};

% Treg feedback: suppression → tumour killing (dashed red)
\draw[fbarr, tumorcell] (suppr.south) |- ([yshift=3pt]kill.west)
    node[near end, above, font=\footnotesize, tumorcell!70!black]
    {inhibit};

% ── Annotations ─────────────────────────────────────────────
\node[font=\footnotesize, gray, anchor=west] at (1.0, 0) {initial event};
\node[font=\footnotesize, gray, anchor=east] at (-1.8,-10) {negative feedback};
\node[font=\footnotesize, gray, anchor=west] at (6.0,-10) {effector response};

\end{tikzpicture}
}
\caption{BPMN-style representation of the immune response cascade. Sequential processing from tumor detection through antigen presentation to T~cell activation, with negative feedback from regulatory T~cells (red dashed arrows) that can suppress effector functions at multiple points.}
\label{fig:immune_cascade_bpmn}
\end{figure}

% ============================================================================
\section{Tumor Metabolism}
\label{sec:bg_metabolism}
% ============================================================================

Cancer cells preferentially metabolize glucose through glycolysis even when oxygen is plentiful---the Warburg effect~\cite{liberti2016}.  Although glycolysis yields only 2 ATP per glucose (vs.\ 36 from oxidative phosphorylation), it proceeds faster and feeds biosynthetic pathways for rapid division~\cite{vander_heiden2009}.  Activated T cells undergo similar metabolic reprogramming~\cite{chang2015}, creating competition for the same glucose supply: tumor growth depletes glucose, suppresses immune function, and permits further growth.  This is especially relevant for immunotherapy, because checkpoint inhibitors can only restore T cell function if sufficient metabolic resources remain.  The model captures this through an explicit glucose diffusion field (Chapter~\ref{ch:biological-subsystems}).

% ============================================================================
\section{Treatment Modalities}
\label{sec:bg_treatments}
% ============================================================================

The model implements two drug classes and their combination. Immune checkpoint inhibitors (ICIs) block the PD-1/PD-L1 axis, restoring T cell effector function against tumor cells~\cite{pardoll2012, ribas2018}. Tyrosine kinase inhibitors (TKIs) suppress VEGF-driven angiogenesis, restricting nutrient supply to the tumor while also inducing immunogenic cell death that primes dendritic cells for adaptive immune activation. Combination ICI+TKI therapy---now the standard first-line treatment for advanced RCC following pivotal trials such as KEYNOTE-426~\cite{rini2019} and JAVELIN Renal 101~\cite{motzer2019}---exploits the synergy between these mechanisms: TKI-mediated vascular normalization improves immune cell access to the tumor, while ICI ensures that activated T cells retain cytotoxic capacity. The ARON-1 study~\cite{santoni2024} extends this line of investigation by examining whether the optimal strategy depends on patient sex. Chapter~\ref{ch:treatment} provides the full mechanistic detail of how these treatments are realized in the simulation.

% ============================================================================
\section{Related Work}
\label{sec:related_work}
% ============================================================================

This section positions the present work within the broader computational oncology landscape, covering continuum models, agent-based frameworks, and the direct predecessor of this thesis.

% ----------------------------------------------------------------------------
\subsection{Continuum and PDE-Based Tumor Models}
\label{subsec:rw_continuum}
% ----------------------------------------------------------------------------

Classical mathematical oncology relies on ODE/PDE systems to describe tumor growth as a continuum process. Benzekry et al.~\cite{benzekry2014} showed that standard growth laws (exponential, logistic, Gompertz) can reproduce aggregate growth curves but treat the tumor as a homogeneous mass, discarding intratumoral heterogeneity, spatial gradients, and stochastic cell-level events. Macklin et al.~\cite{macklin2012} advanced the paradigm by coupling reaction-diffusion nutrient transport with tissue mechanics, capturing macroscopic features like necrotic core formation, but their continuum representation still averages over individual cell states. These limitations motivate agent-based approaches whenever cell-level diversity---such as the coexistence of immunogenic and immune-evasive subclones, or stochastic T cell activation and exhaustion---is central to the question under study.

% ----------------------------------------------------------------------------
\subsection{PhysiCell: A Continuous Mechanics-Based Cell Simulator}
\label{subsec:rw_physicell}
% ----------------------------------------------------------------------------

PhysiCell~\cite{ghaffarizadeh2018} represents cells as off-lattice agents governed by continuous mechanics (adhesion, repulsion, deformation) coupled with PDE-solved chemical fields. It excels at biophysical realism---spheroid compaction, migration through extracellular matrix---but its computational cost scales with mechanical solver fidelity, and its architecture is oriented toward biomechanical rather than immunological questions. The present work trades sub-cellular mechanical resolution for the ability to represent thirteen distinct immune cell types, gene-level tumor heterogeneity, and metabolic competition on a shared glucose field.

% ----------------------------------------------------------------------------
\subsection{HAL: Hybrid Automata Library}
\label{subsec:rw_hal}
% ----------------------------------------------------------------------------

The Hybrid Automata Library (HAL)~\cite{bravo2013} is a Java-based toolkit that combines on-lattice agents with continuous PDE layers for diffusible substances, sharing the conceptual approach of coupling discrete cells with a nutrient field. However, HAL models typically aggregate immune cells into a single generic effector population, whereas the present model distinguishes thirteen subtypes with specific behavioral rules. HAL also lacks native parallelism, limiting scalability for the large parameter sweeps required by Bayesian calibration.

% ----------------------------------------------------------------------------
\subsection{Agent-Based Models of Cancer}
\label{subsec:rw_cancer_abms}
% ----------------------------------------------------------------------------

Norton et al.~\cite{norton2019} showed that spatial patterns of immune infiltration determine immunotherapy outcomes in breast cancer.  Wang et al.~\cite{wang2015} demonstrated cross-scale feedback by embedding intracellular signaling ODEs within agents, a philosophy adopted here through the 16-gene DNA system.  Anderson and Quaranta~\cite{anderson2008} showed that spatial structure alters subclone competition.  The present work builds on these foundations while targeting RCC with immune cell diversity exceeding most prior models.

% ----------------------------------------------------------------------------
\subsection{The Merelli et al.\ Mesa-Based RCC Model}
\label{subsec:rw_merelli}
% ----------------------------------------------------------------------------

The most direct predecessor is the RCC ABM by Merelli et al.~\cite{merelli2021}, implemented in Mesa on a 2D lattice with tumor cells and a simplified immune set.  That model reproduced qualitative RCC progression features but lacked metabolic modeling, genetic heterogeneity, and quantitative validation.  The present thesis extends this in five directions: Repast4Py with MPI parallelism; a diffusion-based glucose field; a 16-gene DNA system; sex hormone modulation enabling ARON-1 investigation; and Bayesian optimization-based calibration against clinical endpoints.

% ----------------------------------------------------------------------------
\subsection{Comparative Summary}
\label{subsec:rw_comparison}
% ----------------------------------------------------------------------------

Table~\ref{tab:framework_comparison} summarizes these trade-offs. The present work occupies a distinct position, prioritizing immunological detail, metabolic realism, and genetic heterogeneity within a scalable parallel framework.

\begin{table}[htbp]
\centering
\caption{Comparison of the present model with related frameworks and the predecessor model.}
\label{tab:framework_comparison}
\footnotesize
\begin{adjustbox}{max width=\textwidth}
\begin{tabular}{lllll}
\toprule
\textbf{Dimension} & \textbf{This work} & \textbf{PhysiCell} & \textbf{HAL} & \textbf{Merelli et al.} \\
\midrule
Agent types & 18 (tumor, 13 immune, & Custom cell & On-/off-lattice & Tumor + \\
 & 4 environmental) & phenotypes & agents & few immune types \\
\addlinespace
Spatial model & 3D discrete lattice & Continuous & Lattice + & 2D discrete \\
 & (Repast4Py grid) & off-lattice & off-lattice & lattice (Mesa) \\
\addlinespace
Metabolic layer & Diffusion-based & PDE-based & PDE-based & None \\
 & glucose field & oxygen/nutrients & substances & \\
\addlinespace
Immune detail & 13 subtypes with & Generic effectors, & Typically & Simplified \\
 & specific rules & customizable & aggregated & immune types \\
\addlinespace
Genetic model & 16-gene binary & Phenotype & Not standard & None \\
 & DNA per tumor cell & parameters & & \\
\addlinespace
Treatment & ICI + TKI + & Drug diffusion, & Drug fields & Basic \\
modeling & hormone modulation & pharmacokinetics & & treatment rules \\
\addlinespace
Parallelism & MPI-native & OpenMP threads & Single-process & Single-process \\
\addlinespace
Validation & Bayesian opt.\ + & Problem-specific & Problem-specific & Qualitative \\
approach & ARON-1 clinical data & & & plausibility \\
\bottomrule
\end{tabular}
\end{adjustbox}
\end{table}

No prior ABM of the tumor microenvironment combines explicit glucose metabolism, sex hormone modulation, and calibration against sex-stratified clinical data within a single framework.  The present work takes an early step toward precision oncology by calibrating against sex-stratified ARON-1 data, demonstrating that even a single stratification variable yields clinically meaningful differences in predicted treatment response.

With these theoretical foundations in place, the following chapter presents the overall model architecture and design decisions that translate these concepts into a working simulation.